\documentclass[11pt,letterpaper]{article}
\usepackage[margin=1in]{geometry}
\usepackage{graphicx} % Required for inserting images
\usepackage{tikz}
\usepackage{amsthm,amsmath,amssymb}
\usepackage{algorithm}
\usepackage{algpseudocode}
\usepackage{xcolor}
\usepackage{soul}
\usepackage{enumitem}

\usepackage[hidelinks]{hyperref}
\usepackage{cleveref}

% Theorems
\newtheorem{theorem}{Theorem}[section]
\newtheorem{corollary}[theorem]{Corollary}
\newtheorem{lemma}[theorem]{Lemma}
\newtheorem{claim}[theorem]{Claim}
\theoremstyle{definition}
\newtheorem{definition}[theorem]{Definition}
\newtheorem{remark}[theorem]{Remark}
\newtheorem{observation}[theorem]{Observation}

\crefname{theorem}{Theorem}{Theorems}
\Crefname{theorem}{Theorem}{Theorems}
\crefname{corollary}{Corollary}{Corollaries}
\Crefname{corollary}{Corollary}{Corollaries}
\crefname{lemma}{Lemma}{Lemmas}
\Crefname{lemma}{Lemma}{Lemmas}
\crefname{claim}{Claim}{Claims}
\Crefname{claim}{Claim}{Claims}
\crefname{definition}{Definition}{Definitions}
\Crefname{definition}{Definition}{Definitions}
\crefname{remark}{Remark}{Remarks}
\Crefname{remark}{Remark}{Remarks}
\crefname{observation}{Observation}{Observations}
\Crefname{observation}{Observation}{Observations}

\newcommand{\reals}{\mathbb{R}}
\newcommand{\Span}{\mathrm{span}}
\newcommand{\coreq}{\mathrm{coreq}}

\newcommand{\calA}{\mathcal{A}}
\newcommand{\calB}{\mathcal{B}}
\newcommand{\calC}{\mathcal{C}}
\newcommand{\calD}{\mathcal{D}}
\newcommand{\calE}{\mathcal{E}}
\newcommand{\calF}{\mathcal{F}}
\newcommand{\calG}{\mathcal{G}}
\newcommand{\calH}{\mathcal{H}}
\newcommand{\calI}{\mathcal{I}}
\newcommand{\calJ}{\mathcal{J}}
\newcommand{\calK}{\mathcal{K}}
\newcommand{\calL}{\mathcal{L}}
\newcommand{\calM}{\mathcal{M}}
\newcommand{\calN}{\mathcal{N}}
\newcommand{\calO}{\mathcal{O}}
\newcommand{\calP}{\mathcal{P}}
\newcommand{\calQ}{\mathcal{Q}}
\newcommand{\calR}{\mathcal{R}}
\newcommand{\calS}{\mathcal{S}}
\newcommand{\calT}{\mathcal{T}}
\newcommand{\calU}{\mathcal{U}}
\newcommand{\calV}{\mathcal{V}}
\newcommand{\calW}{\mathcal{W}}
\newcommand{\calX}{\mathcal{X}}
\newcommand{\calY}{\mathcal{Y}}
\newcommand{\calZ}{\mathcal{Z}}

\newcommand{\cost}{\mathrm{COST}}
\newcommand{\ALG}{\mathrm{ALG}}
\newcommand{\OPT}{\mathrm{OPT}}
\newcommand{\Flow}{\textsc{Flow}}
\newcommand{\ouralg}{\textsc{LineTracker}}

\newcommand{\EE}{\mathbb{E}}

\newcommand{\interval}[1]{\lfloor #1 \rceil}
\newcommand{\dis}[1]{#1}
\newcommand{\off}{\mathrm{off}}
\newcommand{\on}{\mathrm{on}}
\newcommand{\org}{\mathrm{org}}
\newcommand{\rel}{\mathrm{rel}}
\newcommand{\new}{\mathrm{new}}

\newcommand{\solution}{\medskip\noindent\textbf{Solution.}\par\medskip}

\newcommand{\Var}{\mathrm{Var}}

\title{21-701: Discrete Mathematics \\ Homework 3}
\author{Poon Tze-Yang}
\date{Due October 5, 2026}


\begin{document}

\maketitle

\noindent Instructor: Wesley Pegden

\begin{enumerate}[series=q]
\item Let $G$ be a finite graph on the vertex set $V$ and for each $v\in V$ let $L(v)\subseteq\mathbb{Z}$ be a finite list of colors. Prove that if there is a $d\geq 1$ such that each $L(v)$ has size at least $10d$ and such that for all $v$ and for all $c\in L(v)$, we have that
\[
|\{u\in V \mid u\sim v \text{ and } c\in L(u)\}| \leq d,
\]
then there is a proper coloring of $G$ (i.e., where adjacent vertices get distinct colors) where each vertex $v$ gets a color from its list $L(v)$.
\end{enumerate}

\solution

For each vertex create a new list of colors $L'(v)$ comprising the first $10d$ colors of $L(v)$.
Uniformly randomly assign each vertex one color from $L'(v)$.
For each pair of vertices $u,v$ such that $u\sim v$ and color $c$ such that $c\in L'(u)\cap L'(v)$, let $A_{u,v,c}$ be the event that vertices $u$ and $v$ are both colored $c$.
Then we have that
\begin{align*}
    \Pr(A_{u,v,c}) &= \frac{1}{|L'(u)|}\cdot \frac{1}{|L'(v)|} \\
    &=\frac{1}{100d^2}.
\end{align*}

Let $\calA_{u,v}$ be the collection of events $A_{x,y,c}$ where $x,y\notin \{u,v\}$.
$A_{u,v,c}$ is a function of the colors assigned to $u$ and $v$, and all events in $\calA_{u,v}$ are functions of the colors assigned to vertices other than $u$ and $v$.
Since each vertex's color is independently assigned, $A_{u,v,c}$ is mutually independent from $\calA_{u,v}$.

By the condition $|\{u\in V \mid u\sim v \text{ and } c\in L(u)\}| \leq d$, there are at most $|L'(v)|d = 10d^2$ ways to select a color $c\in L'(v)$ and an adjacent vertex $w$ where $c\in L'(w)$. 
Similarly, there are at most $10d^2$ ways to do the same for $u$. 
Each of the total $20d^2$ ways to get a monochromatic edge involving $u$ or $v$ above corresponds to one of the events adjacent to $A_{u,v,c}$ in the dependency graph.
Hence, the degree of $A_{u,v,c}$ in the dependency graph is at most $20d^2$.

Therefore, we have that 
\begin{align*}
    e \Pr(A_{u,v,c}) (\text{(degree of }A_{u,v,c}\text{ in the dependency graph)}+1) &\leq \frac{e (20 d^2+1)}{100 d^2} \leq 1
\end{align*}
for $d\geq 1$. Hence, by the symmetric Lovasz Local Lemma, there exists an assignment of colors from each $L'(v)$ to each $v$ such that there is no monochromatic edge.

\begin{enumerate}[resume=q]
\item Prove that for any $\varepsilon>0$, there is an $N_\varepsilon\in\mathbb{N}$ such that for every $T$, there is a binary sequence of length $T$ in which for all $n>N_\varepsilon$, any two identical blocks of length $n$ are separated by distance at least $(2-\varepsilon)^n$.
\end{enumerate}

\solution

Fix $T$ and let $b_1b_2\cdots b_T$ be a sequence of independent uniformly random bits.
For $n>N$ and $1\leq i<j\leq T-n+1$ with $j-i<(2-\varepsilon)^n$, let $A_{i,j,n}$ be the event that $b_{i+t}=b_{j+t}$ for all $0\leq t<n$.
Any sequence in which no $A_{i,j,n}$ occurs has the required property.

\paragraph{Probability.}
Conditioned on $b_1,\dots,b_{j+t-1}$, the bit $b_{j+t}$ equals the earlier bit $b_{i+t}$ with probability $1/2$.
Hence $\Pr(A_{i,j,n})=2^{-n}$, even when the two blocks overlap.

\paragraph{Dependency graph.}
$A_{i,j,n}$ is determined by the at most $2n$ bits in $S=[i,i+n-1]\cup[j,j+n-1]$, so it is mutually independent of all events whose bits avoid $S$.
Join two events if their sets of bits intersect.
A neighbour $A_{i',j',m}$ of $A_{i,j,n}$ is specified by a position $k\in S$ (at most $2n$ ways), which of its two blocks contains $k$ ($2$ ways), the offset of $k$ in that block ($m$ ways), and the start of its other block, which is at distance less than $(2-\varepsilon)^m$ (fewer than $(2-\varepsilon)^m$ ways).
So $A_{i,j,n}$ has at most $4nm(2-\varepsilon)^m$ neighbours of block length $m$.

\paragraph{Local Lemma.}
Set $x_{i,j,n}=x_n=\frac{1}{(2-\varepsilon)^n n^3}\leq\frac12$.
We use $1-x\geq e^{-2x}$ for $0\leq x\leq\frac12$ (the convex function $e^{-2x}$ lies below its chord from $(0,1)$ to $(\frac12,e^{-1})$, which lies below $1-x$), and $\sum_{m>N}\frac{1}{m^2}\leq\sum_{m>N}\frac{1}{m(m-1)}=\sum_{m>N}\left(\frac{1}{m-1}-\frac{1}{m}\right)=\frac1N$.
For every event $A=A_{i,j,n}$,
\begin{align*}
    x_n\prod_{B\sim A}(1-x_B)
    &\geq x_n\prod_{m>N}(1-x_m)^{4nm(2-\varepsilon)^m} \\
    &\geq x_n\prod_{m>N}(\exp(\frac{1}{(2-\varepsilon)^mm^3}))^{4nm(2-\varepsilon)^m} \\
    &\geq x_n\exp\Big(-8n\sum_{m>N}\frac{1}{m^2}\Big) \\
    &\geq \frac{e^{-8n/N}}{(2-\varepsilon)^n n^3} \\
    &=\left((2-\varepsilon)e^\frac{8}{N}\right)^{-n} 
\end{align*}
For a sufficiently large $N$ which depends only on $\varepsilon$, the term in the base is smaller than $2$, so in particular we have that the whole term is larger than $2^{-n}=\Pr(A)$.
By the asymmetric Lov\'asz Local Lemma, with positive probability no $A_{i,j,n}$ occurs, so the required sequence exists.


\begin{enumerate}[resume=q]
\item Prove that the following is true for sufficiently large $k$:

For any subset $S$ of $\mathbb{N}$ of size $k$ and any natural number $N$, there is a $3$-coloring of the set $[N] = \{1,\dots,N\}$ so that every translate of $S$ that is a subset of $[N]$ contains a natural number of every color.
\end{enumerate}

\solution

For the given $S$, let $S_i$ be the translation of $S$ by $i$, so every number $n$ in $S$ gets set to $(n-1+i)\% N +1$.
Assign colors to each number independently, uniformly at random.
Let $A_{i,c}$ for some color $c\in [3]$ be the event that $S_i$ does not contain the color $c$.
Then for all $A_{i,c}$,
\begin{align*}
    \Pr(A_{i,c}) = \left(\frac{2}{3}\right)^k =: P.
\end{align*}

$A_{i,c}$ is determined by the colorings of the $k$ vertices in $S_i$, so it is mutually independent of all events which avoid the vertices of $S_i$.
In the dependency graph, join two events $A_{i,c}$ and $A_{j,c'}$ if $S_i\cap S_j \neq \emptyset$.
Each neighbor of $A_{i,c}$ overlaps at least one vertex of $S_i$ ($k$ ways), and that vertex takes some position in $S_j$ ($k$ ways), and $c'$ is some color in $[3]$ (3 ways).
Hence, $A_{i,c}$ has at most $3k^2$ neighbors.

We have that, 
\begin{align*}
    e \Pr(A_{i,c}) (3k^2+1) = e \left(\frac{2}{3}\right)^k (3k^2+1)<1.
\end{align*}
for sufficently large $k$, so by the Lovasz Local Lemma, there is at least one coloring that avoids all non-rainbow $S_i$'s.

\begin{enumerate}[resume=q]
\item Let $\mathrm{R}(k,k)$ denote the least $n$ such that every $2$-coloring of the edges of $K_n$ contains a monochromatic $K_k$. Use the Local Lemma to prove that
\[
\mathrm{R}(k,k) \geq (1+o(1))\frac{\sqrt{2}}{e}\,k\,2^{k/2}.
\]
\end{enumerate}

\solution

Assign colors to edges of $K_n$ independently, uniformly at random. 
For each subset of $k$ vertices $S$ and color $c\in[2]$, let $A_{S,c}$ be the event that $S$ is a $K_k$ clique of color $c$.
We have that this event occurs exactly if all edges of $S$ are colored $c$, so
\begin{align*}
    \Pr(A_{S,c}) = \left(\frac{1}{2}\right)^{\binom{k}{2}} = \left(\frac{1}{2}\right)^{\frac{k(k-1)}{2}} .
\end{align*}

$A_{S,c}$ is determined by the colorings of the edges of the clique on $S$, so it is mutually independent of the events $A_{S',c'}$ where $|S\cap S'|\leq 1$, since such sets cannot share edges.
For each $S$, there are $\binom{k}{\ell}\binom{n-\ell}{k-\ell}$ sets $S'$ of vertices of size $k$ which overlap with $S$ on $\ell$ vertices.
There are $2$ colors for each such $S'$ which correspond to events $A_{S',c'}$.
Hence, letting the degree of $A_{S',c}$ in the dependency graph be $D$,
\begin{align*}
    D &\leq 2 \sum_{\ell=2}^k \binom{k}{\ell}\binom{n-k}{k-\ell} \\
    &\leq 2\sum_{\ell=2}^k \binom{k}{\ell}\frac{n^{k-\ell}}{(k-\ell)!}.
\end{align*}
For large $n$, the $\ell=2$ term dominates, so 
\begin{align*}
    D \leq 2\cdot \frac{k(k-1)}{2}\cdot \frac{n^{k-2}}{(k-2)!} + o(n^{k-2}) \leq k^2\left(\frac{en}{k-2}\right)^{k-2}
\end{align*}
For the Local Lemma to hold, we need
\begin{align*}
    eP(D+1)<1 \Rightarrow D< \frac{1}{e}\cdot 2^{\frac{k(k-1)}{2}}
\end{align*}
The following condition suffices:
\begin{align*}
    &k^2\left(\frac{en}{k-2}\right)^{k-2} \leq \frac{1}{e}\cdot 2^{\frac{k(k-1)}{2}} \\
    \Rightarrow & \frac{en}{k-2}\leq \left(\frac{1}{ek^2}\right)^{\frac{1}{k-2}}\cdot 2^{\frac{k(k-1)}{2(k-2)}} \\
    \Rightarrow &n\leq \left(\frac{1}{ek^2}\right)^{\frac{1}{k-2}}\cdot 2^{\frac{k(k-1)}{2(k-2)}}\cdot \frac{k-2}{e}.
\end{align*}
Note that when $k\to \infty$, the above term tends to 
\begin{align*}
    \frac{k}{2}2^\frac{k}{2}.
\end{align*}
Hence, when $n$ is smaller than this quantity, there exists a 2-coloring of the edges that gives a monochromatic $K_k$. 
Thus $R(k,k)$ is lower bounded by this quantity.


\begin{enumerate}[resume=q]
\item Define $\mathrm{R}^c(k_1,\dots,k_c)$ to be the smallest $n$ so that any $[c]$-coloring of the edges of $K_n$ has, for some $i$, a $K_{k_i}$ subgraph that is all color $i$. State and prove an upper bound on $\mathrm{R}^c(3,\dots,3)$.
\end{enumerate}

\solution


\qed


The goal of the final problem is to see a combinatorial route to going beyond the probabilistic method to prove lower bounds for hypergraph Ramsey numbers. Specifically, we will be proving
\[
\mathrm{R}_3(k;4) \geq 2^{\mathrm{R}_2(k-1;2)-1}
\]
where $\mathrm{R}_r(k;c)$ is the minimum $n$ such that in any $c$-coloring of the hyperedges of the $r$-uniform hypergraph on $[n]$, there is a monochromatic complete subhypergraph on $k$ vertices. Note that $\mathrm{R}_2(k-1;2)$ is the familiar $2$-color graph Ramsey number $\mathrm{R}(k-1)$ and has a lower bound of order $2^{(k-1)/2}$, so we will be showing a lower bound that is a tower of height $2$. So in this example you are using the $2$-color graph Ramsey number to give a lower bound on the $4$-color $3$-uniform hypergraph Ramsey number. (Versions of the argument below also exist for higher uniformities without needing to increase the number of colors.)

Let $n = \mathrm{R}_2(k-1;2)-1$ (note that for this $n$ there is a coloring of the edges of $K_n$ without a monochromatic $K_{k-1}$ subgraph).

Now we let $V$ be the set of binary strings
\[
b_1b_2b_3\cdots b_n, \qquad b_i\in\{0,1\},
\]
so that $N = |V| = 2^n$. Our goal is to $4$-color all three-tuples of $V$ so that there is no complete monochromatic subhypergraph on $k$ vertices of $V$, assuming we are given a two-coloring $c:\binom{[n]}{2}\to\{\text{red},\text{blue}\}$ of the pairs in $[n]$ without a monochromatic $K_{k-1}$ subgraph.

Let $<$ be the lexicographic order on the strings $\sigma\in V$ (i.e., we order the strings as they would be ordered if interpreted as binary representations of natural numbers).

Now for a given triple $\sigma_1,\sigma_2,\sigma_3$, assume without loss of generality that $\sigma_1<\sigma_2<\sigma_3$, and let $i\in[n]$ be the first digit in which $\sigma_1$ and $\sigma_2$ differ, and let $j\in[n]$ be the first digit in which $\sigma_2$ and $\sigma_3$ differ.

We give the triple $\{\sigma_1,\sigma_2,\sigma_3\}$ the color $1$, $2$, $3$, or $4$ according to which of the following cases holds:
\begin{enumerate}[label=\arabic*.]
\item If $i<j$ and in the given coloring $c$ of $\binom{[n]}{2}$, $\{i,j\}$ is red,
\item If $i>j$ and in the given coloring $c$ of $\binom{[n]}{2}$, $\{i,j\}$ is red,
\item If $i<j$ and in the given coloring $c$ of $\binom{[n]}{2}$, $\{i,j\}$ is blue,
\item If $i>j$ and in the given coloring $c$ of $\binom{[n]}{2}$, $\{i,j\}$ is blue.
\end{enumerate}

\begin{enumerate}[resume=q]
\item Prove that there is no monochromatic complete ($3$-uniform) subhypergraph on $k$ vertices, with this coloring.
\end{enumerate}

\solution


\end{document}
